% year: תשע"ח
% type: מבחן
% semester: א
% moed: ב
% number: 3ב
% points: 10
% categories: ivt, mvt
% source: exams/מבחנים/תשעח/מועד ב תשעח.pdf

\begin{question}
הוכיחו כי יש לפולינום $x^3-4x^2+7x+13$ שורש ממשי אחד ויחיד.
\end{question}

\begin{hint}
קיום — משפט ערך הביניים. יחידות — הראו שהנגזרת חיובית תמיד.
\end{hint}

\begin{solution}
יהי $p(x)=x^3-4x^2+7x+13$.

\textbf{קיום:} $p$ רציף (פולינום). $p(-2)=-8-16-14+13=-25<0$ ו-$p(0)=13>0$. לפי משפט ערך הביניים קיים $c\in(-2,0)$ עם $p(c)=0$.

\textbf{יחידות:} $p'(x)=3x^2-8x+7$. הדיסקרימיננטה היא $64-84=-20<0$ והמקדם המוביל חיובי, ולכן $p'(x)>0$ לכל $x\in\R$. מכאן $p$ עולה ממש על $\R$ (מסקנה ממשפט לגרנז'), ולכן אינה יכולה לקבל את הערך $0$ פעמיים.

(לחלופין: אם היו שני שורשים $a<b$, לפי משפט רול היה $c\in(a,b)$ עם $p'(c)=0$, סתירה.)

לכן לפולינום שורש ממשי אחד ויחיד (הנמצא בקטע $(-2,0)$; בקירוב $c\approx-1.054$). $\blacksquare$
\end{solution}
